\documentclass[12pt,a4paper]{article}
\usepackage{ugmath}
\usepackage{float}
\usepackage{placeins}
\usepackage{array}
\usepackage[labelfont=bf,labelsep=space]{caption}
\newcommand{\paren}[1]{\left( #1 \right)}
\newcommand{\alumno}{Ricardo León Martínez}
\newcommand{\materia}{Elementos de Probabilidad y Estadistica}
\newcommand{\profesor}{Claudia Reynoso Alcántara}
\newcommand{\tarea}{Tarea 11}
\newcommand{\fecha}{5/6/2026}

\begin{document}

    \begin{center}
        {\large\textbf{UNIVERSIDAD DE GUANAJUATO}}\\[0.3cm]
        {\normalsize\textbf{DIVISIÓN DE CIENCIAS NATURALES Y EXACTAS}}\\
        {\normalsize\textbf{CAMPUS GUANAJUATO}}\\[1cm]

        {\Large\textbf{\tarea\ (\materia)}}\\[1cm]
    \end{center}

    \noindent
    \textbf{Nombre:} \alumno 
    \hfill 
    \textbf{Fecha:} \fecha 
    \hfill 
    \textbf{Calificación:} \rule{3cm}{0.4pt}

    \vspace{0.3cm}

    \section*{Parte I. Valor esperado y varianza}

    \begin{mdframed}[style=mdpinkbox,frametitle={Ejercicio 1}]
        Sea $X$ una variable aleatoria discreta tal que $P(X=a)=p$ y $P(X=b)=1-p$ para algunos $a,b,p\in\R$ con $a\leq b$
        y $a\lt p\lt1$.
        \begin{enumerate}
            \item[(a)] Encuentre una expresión para $\E[X^{n}]$ para cada $n\geq1$ y determine $\Var(X)$.
            \item[(b)] Suponiendo que $a=1$ y $b=-1$, encuentre $c\neq1$ tal que $\E[c^{X}]=1$. 
        \end{enumerate}
    \end{mdframed}

    \begin{mdframed}[style=mdpinkbox,frametitle={Ejercicio 2}]
        Sea $X$ una variable aleatoria con valor esperado $\mu=\E[X]$ y varianza $\sigma^{2}=\Var(X)$.
        \begin{enumerate}
            \item[(a)] Encuentre el valor esperado y varianza de $Y=(X-\mu)/\sigma^{2}$.
            \item[(b)] Sean $a,b\in\R$. Encuentre el valor esperado y varianza de $Z=a+bY$.
        \end{enumerate}
    \end{mdframed}

    \section*{Parte II. Distribución exponencial}

    Recordemos que una variable aleatoria $X$ tiene \textbf{distribución exponencial} de parámetro $\lambda\gt0$ si tiene densidad
    \begin{equation*}
        f(r)=
        \begin{cases}
            xe^{-\lambda r},\hfill\quad\text{si }r\geq0,\\
            0,\hfill\quad\text{en cualquier otro caso}.
        \end{cases}
    \end{equation*}

    \begin{mdframed}[style=mdpinkbox,frametitle={Ejercicio 3}]
        Sean $X$ una variable aleatoria exponencial de parámetro $\lambda\gt0$. Demuestra que para todo entero $n\geq1$ se cumple que
        \begin{equation*}
            \E[X^{n}]=\frac{n}{\lambda}\cdot\E[X^{n-1}].
        \end{equation*}
        Encuentre $E[X]$ y $E[X^{2}]$ y concluya que $\Var(X)=\frac{1}{\lambda^{2}}$. \textbf{Sugerencia:} Use integración por partes
        con $u=r^{n}$ y $dv=\lambda e^{-\lambda r}dr$.
    \end{mdframed}

    \section*{Parte III. Distribución gamma}

    La \textbf{Distribución gamma} $\Gamma:\R_{\gt0}\to\R$ esta dada por
    \begin{equation*}
        \Gamma(\alpha)=\int_{0}^{\infty}e^{-r}r^{\alpha-1}dr\quad\text{para todo}\quad\alpha\gt0.
    \end{equation*}

    \begin{mdframed}[style=mdpinkbox,frametitle={Ejercicio 4}]
        Demuestra usando integración por partes que $\Gamma(\alpha)=(\alpha-1)\Gamma(\alpha-1)$ para $\alpha\gt1$. Encuentre
        $\Gamma(1)$ y concluya que $\Gamma(n)=(n-1)!$.
    \end{mdframed}

    Recordemos que una variable aleatoria $X$ tiene \text{distribución gamma} de parametros $(\lambda,\alpha)$ con $\lambda\gt0$
    y $\alpha\gt0$ si tiene densidad
    \begin{equation*}
        f(r)=
        \begin{cases}
            \frac{\lambda e^{-\lambda r}(\lambda r)^{\alpha-1}}{\Gamma(\alpha)},\hfill\text{si } r\geq0,\\
            0,\qquad\text{en cualquier otro caso}.
        \end{cases}
    \end{equation*}

    \begin{mdframed}[style=mdpinkbox,frametitle={Ejercicio 5}]
        Sea $X$ una variable aleatoria gamma de parámetros $(\lambda,\alpha)$ con $\lambda\gt0$ y $\alpha\gt0$. Encuentre
        $\E[X]$ y $\E[X^{2}]$ y concluya que $\Var(X)=\alpha/\lambda^{2}$. \textbf{Sugerencia.} Escriba $\lambda r^{n}e^{-\lambda r}(\lambda r)^{n+\alpha-1}$
        y considere el cambia de variable $s=\lambda r$.
    \end{mdframed}

    \section*{Parte IV. Distribución beta}

    Recordemos que variable aleatoria $X$ tiene \textbf{distribución beta} de parámetros $(a,b)$ con $a,b\gt0$ si tiene densidad
    \begin{equation*}
        f(r)=
        \begin{cases}
            \frac{1}{C(a,b)}r^{a-1}(1-r)^{b-1},\quad\text{ si }0\lt r\lt 1,\\
            0,\hfill\text{en cualquier otro caso},
        \end{cases}
    \end{equation*}
    donde
    \begin{equation*}
        C(a,b)=\int_{0}^{1}r^{a-1}(1-r)^{b-1}dr.
    \end{equation*}

    \begin{mdframed}[style=mdpinkbox,frametitle={Ejercicio 6}]
        Sea $X$ una variable aleatoria beta de parametros $(a,b)$ con $a,b\gt0$. Demuestre que
        \begin{equation*}
            \E[X]=\frac{a}{a+b},\quad\text{ y }\quad\Var(X)=\frac{ab}{(a+b)^{2}(a+b+1)}.
        \end{equation*}
        \textbf{Sugerencia.} Use el hecho de que la función  gamma satisface que $\frac{\Gamma(a)\Gamma(b)}{\Gamma(a+b)}=C(a,b)$.
    \end{mdframed}

    \section*{Parte V. Integral de Stieltjes}

    Sea $X$ una variable aleatoria y sea $F_{x}:\R\to\R$ su función de distribución. Recordemos que si una función $g:\R\to\R$ es
    stieltjes integrable con respecto a $F_{x}$ en el intervalo $[a,b]$ con $a\leq b$, entonces
    \begin{equation*}
        \int_{a}^{b}g(r)dF_{X}(r)=\lim_{n\to\infty}\sum_{i=1}^{n}g(r_{i})(F_{x}(r_{i})-F_{x}(r_{i-1}))
    \end{equation*}
    donde para cada entero $n\geq1$ tomamos
    \begin{equation*}
        r_{0}=a,\quad r_{1}=a+\frac{b-a}{n},\quad r_{2}=a+2\frac{b-a}{n},\quad r_{3}=a+3\frac{b-a}{n},\cdots,r_{n-1}=a+(n-1)\frac{b-a}{n},\quad r_{n}=b.
    \end{equation*}
    Más aún, si
    \begin{equation*}
        \int_{-\infty}^{+\infty}g(r)dF_{X}(r):=\lim_{a\to+\infty}\int_{-a}^{+a}g(r)dF_{X}(r)
    \end{equation*}
    existe, lo llamamos el valor esperado de $g(X)$ y lo denotamos por $\E[g(X)]$. En particular, el valor esperado de $X$ es
    \begin{equation*}
        \E[X]=\int_{-\infty}^{+\infty}rdF_{X}(r):=\lim_{a\to+\infty}\int_{-a}^{+a}rdF_{X}(r)
    \end{equation*}

    \begin{mdframed}[style=mdpinkbox,frametitle={Ejercicio 7}]
        Sea $X:\Omega\to\R$ una variable aleatoria no negativa, es decir, $X(\omega)\geq0$ para todo $\omega\in\Omega$.
        Supongamos que $X$ tiene valor esperado $\E[X]$. Demuestre que $\E[X]\geq0$. \textbf{Sugerencia.} Demuestre que
        $F_{X}(r)=0$ si $r\lt0$, y por ende que $\int_{-a}^{+a}rdF_{X}(r)=\int_{0}^{+a}rdF_{X}(r)$.
    \end{mdframed}

\end{document}